% !PRINTED 255
\weberpage{255}
%<<CONT
hörenden Transformationsgleichung durch die für den Grad $p^2$
rational ausdrücken und diesen Ausdruck in die zu $p$ gehörige
Transformationsgleichung einsetzen.

Nach 3., Formel (7) wird die Gleichung
\weq{15}{x^4 + 18x^2 + \gamma_3(\omega)x - 27 = 0}
von den beiden Funktionen
\weq{16}{-\left(\frac{\eta\left(\frac{\omega}{3}\right)}{\eta(\omega)}\right)^6, \quad 27\left(\frac{\eta(3\omega)}{\eta(\omega)}\right)^6}
befriedigt. Bezeichnen wir den ersten dieser Werte mit $x$, so geht
der zweite durch die Substitution $\left(\omega,\frac{\omega}{3}\right)$
in $\dfrac{-27}{x}$ über, und folglich wird durch $x$ auch die
Gleichung
\weq{17}{27^3 + 18\cdot27x^2 - \gamma_3\left(\frac{\omega}{3}\right)x^3 - x^4 = 0}
befriedigt. Setzen wir nun
\weq{18}{y = \left(\frac{\eta\left(\frac{\omega}{9}\right)}{\eta(\omega)}\right)^3,}
so geht $x$ durch die Substitution $\left(\omega,\frac{\omega}{3}\right)$
in $\dfrac{y^2}{x}$ über, so daß man auch die Gleichung erhält:
\weq{19}{y^8 + 18y^4x^2 + \gamma_3\left(\frac{\omega}{3}\right)y^2x^3 - 27x^4 = 0,}
und durch Elimination von $\gamma_3\left(\frac{\omega}{3}\right)$ aus
(17) und (19) erhält man nach Weghebung des Faktors $y^2+27$
\weq{20}{x^4 - 18x^2y^2 - y^2(y^4 - 27y^2 + 27^2) = 0.}

Löst man diese quadratische Gleichung nach $x^2$ auf, so folgt:
\weq{21}{x^2 = y^3 + 9y^2 + 27y = (y+3)^3 - 27,}
wo über das Zeichen durch Einsetzen der Anfangsglieder der
Entwickelungen entschieden wird, am einfachsten wohl, da diese
Gleichung (nach § 69) auch für
\[
x = 27\left(\frac{\eta(3\omega)}{\eta(\omega)}\right)^6, \quad y = 27\left(\frac{\eta(9\omega)}{\eta(\omega)}\right)^3
\]
erfüllt wird, indem man
\[
x = 27q + \cdots, \quad y = 27q^2 + \cdots
\]
setzt.
